\section{Savage and Naked We Come into the World}

\begin{quotex}
We have to think that in an age where the spoken and printed word holds such a considerable place, it is important that the greatest number of Frenchmen know the meaning of the words they use: that is not possible without Greek and Latin. We must realize that at a time where professional technocrats push the spirit of men more and more into the most narrow specialties, a wide-ranging general culture should be the most widespread possible or the vast majority of people will no longer know how to speak or think: that culture is no longer possible without Latin or Greek.
\flright{\textsc{Charles Maurras}\footnote{My translation from \url{http://maurras.net/2011/07/07/souvenirs-de-lettre/}}}
\end{quotex}


If that was true of the 1920s, it is immeasurably truer today. Besides the printed and spoken word, we now have the electronic word which is measured in terabytes or petabytes, that is, solely in terms of quantity. Yet the capacity of the human mind has not changed, so most of that is pure repetition or else wasted. Despite this vast quantity of encoded words, a ``wide-ranging general culture'' is even more elusive than ever.

Quite to the contrary, there seems to be a concerted effort to prevent any widespread dissemination of High Culture. A people without a culture is ``free'' to choose any culture it wants. However, this doesn't take into account that the dissemination of culture is dominated by the few with the resources, talent, and ability to dominate the available attention span of the human mind in a way reminiscent of the old joke: ``When we want your opinion, we'll tell it to you!''

Without knowledge of the sources of ideas from the Greeks and Romans, the majority of people no longer know how to speak or think. Whatever thought that cannot fit on a placard is not worth thinking. This is praised as the democratic ideal, since such thinking is available to anyone. Charles Maurras, in The Battle of the Humanities\footnote{\url{http://maurras.net/textes/117.html}}, explains what is lost in such an attitude.

\begin{quotex}
The accumulated capital of the human race doesn't come to new-borns either as a whole or in equally sized parts. Innumerable streams, of infinite variety, make everything flow in doses not less diverse than what our predecessors did or dreamed before us. The value of all personal effort is dominated by this immense historic ratio in virtue of which, according to the beautiful saying of \textbf{Auguste Comte}, the living are more and more necessarily governed by the dead, and as \textbf{Maurice Barrès} said, everyone alive, by his own particular dead. This beneficial necessity is the source of civilization. But for quite some time democracy has risen up against that condition of a civilized order; it has chosen barbarism that can totally begin anew with each individual who comes, savage and naked, into the world. It is to this humanity of the caves that democracy wants to lead us back.
\end{quotex}


The debate between nature and nurture remains moot. Many want to believe that genetics is all that matters, but, without gainsaying its power, that is to remain at the level of Destiny. Yet, neither is man a blank slate, waiting to be created in the image of some philosopher's futuristic vision. Rather, we depend on the accumulated wisdom of our past to be most compatible with our inherent natures.

If we are not to be dominated by alien ideas whose legitimacy depends on the eradication of an objective knowledge of history, we need to draw on our human capital. As Maurras points out, it comes to us in streams but, as the saying goes, we can't be made to drink. That requires personal effort, or Will, to quench our thirst from these sweet waters rather than from the artificially sweetened beverages of modernity.

We learn from the pagans how to revere antiquity and to continue in the traditions of our ancestors. This is to rethink the thoughts of well-bred men. Even atheists like Maurras and Comte understood this since they understood history. Unlike the majority of intellectuals, they strove to continue the old ways on a new basis. The maintenance of the civilized order is the common goal. There should be no mystery by this point. We have had enough prophets to show us the way out of the cave and into enlightenment.


\flrightit{Posted on 2011-07-07 by Cologero}

\begin{center}* * *\end{center}

\begin{footnotesize}
\begin{sffamily}
\texttt{Will on 2011-07-08 at 00:40 said:}

The destruction of Classical education has indeed been a disaster for the West. Unfortunately, most of the remaining classicists still teaching in the universities are like poor fools sitting on a buried treasure, unaware of what is right underneath them. They concern themselves with textual minutiae and trivialities of grammar, losing sight of the fact that what they hold in their hands is in fact the wisdom tradition of the Western world. It's no wonder so many turn to various foreign or false spiritualities when the very people who ought to be safeguarding and passing down our own wisdom -- Homer, the pre-Socratics, Plato, the New Testament, the Stoics, etc. -- are somehow immune to its life-transforming energies, and do not exemplify its values. Ask a classics professor why one ought to study Greek and Latin and you'll likely hear something like, ``Well, Ted Turner studied classics and now look at him, he's rich!''

Instead of wielding the sword of the Classical tradition against the inferior and destructive thinking of cultural marxism, they hide in their ever-shrinking departments, publishing obscure articles about the use of the genitive absolute in Herodotus, while their students are miseducated by tenured burnouts.

\hfill

\texttt{logres on 2011-07-08 at 08:26 said:}

Even a little Latin can go a long way, as it helps transmute the whole concept of education.

\hfill

\texttt{Izak on 2011-07-08 at 13:37 said:}

Will: Your comment is pretty much 100\% accurate about the professors who are actually decent enough to know and understand classical literature! But the ones who are writing obscure articles on the genitive absolute in such-and-such are part of a dying breed, and they're unfortunately quite good compared to whom they're being replaced by: a slew of lazy, confused, egomaniacal children who feel compelled to rely on ``theory'' to make up for their lack of wisdom. Academia is a business and meaningless novelty is its bread and butter. It will be somewhat amusing watching it crash and burn, although ideally I won't even care enough to notice when it happens.


\hfill

\end{sffamily}
\end{footnotesize}

